\section{Harmonic Analysis and Fourier Transforms}

\subsection{Fourier Transform and Inversion}

\begin{definition}[Fourier Transform]
Let $f \in L^1(\R)$ be a complex-valued, Lebesgue-integrable function. The Fourier transform of $f$, denoted $\hat{f}$, is defined as the function $\hat{f}: \R \to \C$ given by:
\begin{equation}
\hat{f}(\omega) = \int_\R f(t) e^{-i\omega t} \dd t
\end{equation}

\end{definition}

\begin{remark}[Why $L^1(\R)$?]
The requirement that $f \in L^1(\R)$ is strictly necessary to guarantee that the integral converges absolutely for every $\omega \in \R$. Because the complex exponential is bounded by unity ($|e^{-i\omega t}| = 1$), we have:
\begin{equation}
\int_\R |f(t) e^{-i\omega t}| \dd t = \int_\R |f(t)| \dd t = \|f\|_{L^1} < \infty
\end{equation}
Absolute convergence ensures that $\hat{f}(\omega)$ is a well-defined, finite complex number for all $\omega$.
\end{remark}

\begin{lemma}[Riemann--Lebesgue Lemma]
If $f \in L^1(\R)$, then its Fourier transform $\hat{f}$ is continuous, bounded, and vanishes at infinity:
\[
\lim_{|\omega| \to \infty} \hat{f}(\omega) = 0.
\]

\end{lemma}

\begin{proof}
We prove the three properties sequentially.

\textit{1. Boundedness:}
By the triangle inequality for integrals:
\begin{equation}
\begin{aligned}
|\hat{f}(\omega)| &= \left| \int_\R f(t) e^{-i\omega t} \dd t \right| \\
&\leq \int_\R |f(t)| |e^{-i\omega t}| \dd t \\
&= \int_\R |f(t)| \dd t \\
&= \|f\|_{L^1}
\end{aligned}
\end{equation}
Thus, $\hat{f}$ is bounded by the $L^1$ norm of $f$.

\textit{2. Continuity:}
Let $\omega_n$ be a sequence converging to $\omega$. The integrand $f(t) e^{-i\omega_n t}$ converges pointwise to $f(t) e^{-i\omega t}$ for every $t$. Furthermore, the absolute value of the integrand is bounded by $|f(t)|$, which is an integrable dominating function since $f \in L^1(\R)$. By the \textit{Dominated Convergence Theorem}, we can pass the limit inside the integral:
\begin{equation}
\begin{aligned}
    \lim_{n \to \infty} \hat{f}(\omega_n)
    &= \int_\R \lim_{n \to \infty} \left( f(t) e^{-i\omega_n t} \right) \dd t \\
    &= \int_\R f(t) e^{-i\omega t} \dd t = \hat{f}(\omega)
    \end{aligned}
\end{equation}
This proves $\hat{f}$ is continuous.

\textit{3. Decay at Infinity:}
We use a phase-shift trick. Let $\omega \neq 0$. We perform the change of variables $s = t + \pi/\omega$, which implies $t = s - \pi/\omega$ and $\dd t = \dd s$. Notice that
\[
\begin{aligned}
e^{-i\omega(s - \pi/\omega)} &= e^{-i\omega s} e^{i\pi} \\
&= -e^{-i\omega s}.
\end{aligned}
\]
Substituting this into the definition:
\begin{equation}
\begin{aligned}
\hat{f}(\omega) &= \int_\R f(s - \pi/\omega) e^{-i\omega (s - \pi/\omega)} \dd s \\
&= -\int_\R f(s - \pi/\omega) e^{-i\omega s} \dd s
\end{aligned}
\end{equation}

Adding this expression to the original definition of $\hat{f}(\omega)$ (and renaming the dummy variable $s$ back to $t$):
\begin{equation}
2\hat{f}(\omega) = \int_\R \left[ f(t) - f(t - \pi/\omega) \right] e^{-i\omega t} \dd t
\end{equation}
Taking the absolute value and applying the triangle inequality:
\begin{equation}
\begin{aligned}
2|\hat{f}(\omega)| &\leq \int_\R \left| f(t) - f(t - \pi/\omega) \right| \dd t \\
&= \left\| f(\cdot) - f(\cdot - \pi/\omega) \right\|_{L^1}
\end{aligned}
\end{equation}

To show this goes to zero as $|\omega| \to \infty$, we must explicitly prove that translation is continuous in the $L^1$ norm. Let $\tau = \pi/\omega$. We want to show \[
\lim_{\tau \to 0} \|f - f_\tau\|_{L^1} = 0,
\] where $f_\tau(t) = f(t-\tau)$.
Let $\varepsilon > 0$. The space of continuous functions with compact support, $C_c(\R)$, is dense in $L^1(\R)$. Thus, there exists a function $g \in C_c(\R)$ such that $\|f - g\|_{L^1} < \varepsilon/3$. By the triangle inequality:
\begin{equation}
\|f - f_\tau\|_{L^1} \leq \|f - g\|_{L^1} + \|g - g_\tau\|_{L^1} + \|g_\tau - f_\tau\|_{L^1}
\end{equation}
Because the Lebesgue measure is translation-invariant,
\[
\|g_\tau - f_\tau\|_{L^1} = \|g - f\|_{L^1} < \varepsilon/3.
\]

Since $g$ is continuous and has compact support, it is uniformly continuous. Let $K$ be a compact set containing the support of $g$ and $g_\tau$ for small $\tau$. For sufficiently small $|\tau|$, $|g(t) - g(t-\tau)| < (\varepsilon/(3 \mu(K)))$ for all $t$. Thus:
\begin{equation}
\|g - g_\tau\|_{L^1} = \int_K |g(t) - g(t-\tau)| \dd t < \mu(K) \frac{\varepsilon}{3 \mu(K)} = \varepsilon/3
\end{equation}

Combining these, $\|f - f_\tau\|_{L^1} < \varepsilon$. Since $\pi/\omega \to 0$ as $|\omega| \to \infty$, we conclude that $\|f - f_{\pi/\omega}\|_{L^1} \to 0$, which forces $\hat{f}(\omega) \to 0$.
\end{proof}

\begin{theorem}[Fourier Inversion]
If $f \in L^1(\R)$ and its Fourier transform $\hat{f} \in L^1(\R)$, then for almost every $t \in \R$:
\begin{equation}
f(t) = \frac{1}{2\pi} \int_\R \hat{f}(\omega) e^{i\omega t} \dd \omega
\end{equation}
If $f$ is continuous, this equality holds for all $t \in \R$.
\end{theorem}

\begin{proof}
The integral \[
\int \hat{f}(\omega) e^{i\omega t} \dd \omega
\] might not converge absolutely if we merely know $\hat{f} \in L^1$ and attempt to substitute the definition of $\hat{f}$ directly, because the resulting double integral may fail Fubini's conditions. To bypass this, we introduce a regularization factor $e^{-\varepsilon \omega^2}$ (Gaussian damping) and take the limit as $\varepsilon \to 0^+$. Define the regularized integral:
\begin{equation}
I_\varepsilon(t) = \frac{1}{2\pi} \int_\R \hat{f}(\omega) e^{-\varepsilon \omega^2} e^{i\omega t} \dd \omega
\end{equation}
The two facts needed for dominated convergence are
\[
\begin{gathered}
\hat{f} \in L^1 \\
|e^{-\varepsilon \omega^2} e^{i\omega t}| \leq 1,
\end{gathered}
\]

Together, they show that the integrand is dominated by $|\hat{f}(\omega)| \in L^1$. By the Dominated Convergence Theorem, as $\varepsilon \to 0^+$:
\begin{equation}
\lim_{\varepsilon \to 0^+} I_\varepsilon(t) = \frac{1}{2\pi} \int_\R \hat{f}(\omega) e^{i\omega t} \dd \omega
\end{equation}
Now, substitute the definition of \[
\hat{f}(\omega) = \int_\R f(s) e^{-i\omega s} \dd s
\] into $I_\varepsilon(t)$:
\begin{equation}
I_\varepsilon(t) = \frac{1}{2\pi} \int_\R \left[ \int_\R f(s) e^{-i\omega s} \dd s \right] e^{-\varepsilon \omega^2} e^{i\omega t} \dd \omega
\end{equation}

We must justify swapping the order of integration. The absolute value of the integrand is $|f(s)| e^{-\varepsilon \omega^2}$. The double integral of the absolute value is:
\begin{equation}
\begin{aligned}
    \int_\R \int_\R |f(s)| e^{-\varepsilon \omega^2} \dd s \dd \omega
    &= \|f\|_{L^1} \int_\R e^{-\varepsilon \omega^2} \dd \omega \\
    &= \|f\|_{L^1} \sqrt{\frac{\pi}{\varepsilon}} < \infty
    \end{aligned}
\end{equation}

Since the double integral is absolutely convergent, the \textit{Fubini--Tonelli Theorem} permits us to swap the integrals:
\begin{equation}
\begin{aligned}
I_\varepsilon(t)
    &= \int_\R f(s) \underbrace{ \left[ \frac{1}{2\pi} \int_\R e^{-\varepsilon \omega^2} e^{i\omega(t-s)} \dd \omega \right] }_{g_\varepsilon(t-s)} \dd s \\
    &= (f * g_\varepsilon)(t).
\end{aligned}
\end{equation}
Here $*$ denotes convolution. We now explicitly evaluate the inner integral $g_\varepsilon(u)$ by completing the square in the exponent:
\begin{equation}
\begin{aligned}
    -\varepsilon \omega^2 + i\omega u
    &= -\varepsilon \left( \omega^2 - \frac{i\omega u}{\varepsilon} \right) \\
    &= -\varepsilon \left( \left(\omega - \frac{iu}{2\varepsilon}\right)^2 + \frac{u^2}{4\varepsilon^2} \right) \\
    &= -\varepsilon \left(\omega - \frac{iu}{2\varepsilon}\right)^2 - \frac{u^2}{4\varepsilon}
    \end{aligned}
\end{equation}

Substituting $z = \omega - iu/(2\varepsilon)$ and shifting the contour of integration (justified by Cauchy's Integral Theorem, since $e^{-\varepsilon z^2}$ is entire and decays rapidly in horizontal strips):
\begin{equation}
\begin{aligned}
g_\varepsilon(u)
    &= \frac{1}{2\pi} e^{-u^2 / 4\varepsilon} \int_\R e^{-\varepsilon z^2} \dd z \\
    &= \frac{1}{2\pi} e^{-u^2 / 4\varepsilon} \sqrt{\frac{\pi}{\varepsilon}}
    = \frac{1}{\sqrt{4\pi\varepsilon}} e^{-u^2 / 4\varepsilon}
    \end{aligned}
\end{equation}

The family of functions $\{g_\varepsilon\}_{\varepsilon > 0}$ forms a \textit{nascent delta sequence}. We explicitly verify the three required conditions:
\begin{enumerate}
\item \textit{Non-negativity}: $g_\varepsilon(u) > 0$ for all $u \in \R$.
\item \textit{Normalization}: the following identity holds:

\[
\int_\R g_\varepsilon(u) \dd u = 1
\]

(which follows from evaluating the Gaussian integral or noting it is the inverse Fourier transform of $e^{-\varepsilon \omega^2}$ at $t=0$).
\item \textit{Concentration}: for any $\eta > 0$, let $v = u/(4\varepsilon)^{1/2}$. Then:

\begin{equation}
\int_{|u| > \eta} g_\varepsilon(u) \dd u = \frac{2}{\sqrt{\pi}} \int_{|v| > \eta/\sqrt{4\varepsilon}} e^{-v^2} \dd v
\end{equation}
As $\varepsilon \to 0^+$, the lower limit $\eta/(4\varepsilon)^{1/2} \to \infty$, so the integral approaches 0.
\end{enumerate}
By the properties of nascent delta sequences, the convolution $(f * g_\varepsilon)(t)$ converges to $f(t)$ at every point of continuity of $f$. Equating the two limits ($\varepsilon \to 0^+$) yields the inversion formula.
\end{proof}

\subsection{\texorpdfstring{$L^2$}{L2} Theory and Convolution}

\begin{theorem}[Plancherel Theorem]
The Fourier transform extends uniquely to a unitary operator (up to a scaling factor) on $L^2(\R)$. For all $f, g \in L^2(\R)$:
\begin{equation}
\int_\R f(t) \overline{g(t)} \dd t = \frac{1}{2\pi} \int_\R \hat{f}(\omega) \overline{\hat{g}(\omega)} \dd \omega
\end{equation}

\end{theorem}

\begin{proof}
We first prove the identity for $f, g \in \mathcal{S}(\R)$ (the Schwartz space), and then extend it to $L^2(\R)$.

\textit{Step 1: Setup for Schwartz functions.}
Define the time-reversed, complex-conjugated function $\tilde{g}(t) = \overline{g(-t)}$. Since $g \in \mathcal{S}(\R)$, it is clear that $\tilde{g} \in \mathcal{S}(\R)$. Let $h = f * \tilde{g}$. Because $\mathcal{S}(\R)$ is closed under convolution, $h \in \mathcal{S}(\R)$.

\textit{Step 2: Evaluate $h(0)$ in the time domain.}

\begin{equation}
h(0) = (f * \tilde{g})(0) = \int_\R f(t) \tilde{g}(-t) \dd t = \int_\R f(t) \overline{g(t)} \dd t = \langle f, g \rangle_{L^2}
\end{equation}

\textit{Step 3: Evaluate $\hat{h}(\omega)$ in the frequency domain.}
By the Convolution Theorem (proven below),
\[
\hat{h}(\omega) = \hat{f}(\omega) \hat{\tilde{g}}(\omega).
\]
We compute $\hat{\tilde{g}}(\omega)$:
\begin{equation}
\hat{\tilde{g}}(\omega) = \int_\R \overline{g(-t)} e^{-i\omega t} \dd t
\end{equation}
Substitute $s = -t$, so $\dd t = -\dd s$:
\begin{equation}
\begin{aligned}
    \hat{\tilde{g}}(\omega)
    &= \int_\R \overline{g(s)} e^{i\omega s} \dd s \\
    &= \overline{ \left( \int_\R g(s) e^{-i\omega s} \dd s \right) }
    = \overline{\hat{g}(\omega)}
    \end{aligned}
\end{equation}
Thus,
\[
\hat{h}(\omega) = \hat{f}(\omega) \overline{\hat{g}(\omega)}.
\]

\textit{Step 4: Apply the Inversion Theorem.}
Since $h \in \mathcal{S}(\R)$, both $h$ and $\hat{h}$ are in $L^1(\R)$. Applying the Fourier Inversion Theorem at $t=0$:
\begin{equation}
\begin{aligned}
h(0) &= \frac{1}{2\pi} \int_\R \hat{h}(\omega) e^{0} \dd \omega \\
    &= \frac{1}{2\pi} \int_\R \hat{f}(\omega) \overline{\hat{g}(\omega)} \dd \omega \\
    &= \frac{1}{2\pi} \langle \hat{f}, \hat{g} \rangle_{L^2}
    \end{aligned}
\end{equation}
Equating the results from Step 2 and Step 4 yields
\[
\langle f, g \rangle_{L^2} = (1/(2\pi)) \langle \hat{f}, \hat{g} \rangle_{L^2}.
\]

\textit{Step 5: Extension to $L^2(\R)$.}
The Schwartz space $\mathcal{S}(\R)$ is dense in $L^2(\R)$. On this dense subset, the Fourier transform is a linear isometry (up to the $1/(2\pi)^{1/2}$ factor). By the \textit{Bounded Linear Transformation (BLT) Theorem}, any uniformly continuous linear map from a dense subset of a metric space to a complete metric space extends uniquely to the entire space. Since $L^2(\R)$ is complete, the Fourier transform extends uniquely to a bounded linear operator on all of $L^2(\R)$. Because the isometry property holds on the dense subset, it holds on the entire space by the continuity of the norm.
\end{proof}

\begin{corollary}[Parseval's Identity]
Setting $f=g$ in Plancherel's theorem yields the conservation of energy:
\begin{equation}
\int_\R |f(t)|^2 \dd t = \frac{1}{2\pi} \int_\R |\hat{f}(\omega)|^2 \dd \omega
\end{equation}

\end{corollary}

\begin{theorem}[Convolution Theorem]
If $f, g \in L^1(\R)$, then their convolution $f*g \in L^1(\R)$ and the Fourier transform of the convolution is the pointwise product of their Fourier transforms:
\begin{equation}
\widehat{f*g}(\omega) = \hat{f}(\omega) \hat{g}(\omega)
\end{equation}

\end{theorem}

\begin{proof}
By definition,
\[
(f*g)(t) = \int_\R f(\tau) g(t-\tau) \dd \tau.
\]
We evaluate its Fourier transform:
\begin{equation}
\widehat{f*g}(\omega) = \int_\R e^{-i\omega t} \left( \int_\R f(\tau) g(t-\tau) \dd \tau \right) \dd t
\end{equation}
To swap the order of integration, we verify absolute integrability via Fubini--Tonelli:
\begin{equation}
\int_\R \int_\R |f(\tau)| |g(t-\tau)| \dd \tau \dd t = \int_\R |f(\tau)| \left( \int_\R |g(t-\tau)| \dd t \right) \dd \tau = \|f\|_{L^1} \|g\|_{L^1} < \infty
\end{equation}
Since the double integral is absolutely convergent, we swap the integrals:
\begin{equation}
\widehat{f*g}(\omega) = \int_\R f(\tau) \left( \int_\R g(t-\tau) e^{-i\omega t} \dd t \right) \dd \tau
\end{equation}

In the inner integral, we perform the substitution $u = t-\tau$, which implies $t = u+\tau$ and $\dd t = \dd u$. The limits of integration remain $-\infty$ to $\infty$:
\begin{equation}
\begin{aligned}
\int_\R g(u) e^{-i\omega (u+\tau)} \dd u &= e^{-i\omega \tau} \int_\R g(u) e^{-i\omega u} \dd u \\
&= e^{-i\omega \tau} \hat{g}(\omega)
\end{aligned}
\end{equation}
Substituting this back into the outer integral:
\begin{equation}
\begin{aligned}
    \widehat{f*g}(\omega)
    &= \int_\R f(\tau) e^{-i\omega \tau} \hat{g}(\omega) \dd \tau \\
    &= \hat{g}(\omega) \int_\R f(\tau) e^{-i\omega \tau} \dd \tau
    = \hat{g}(\omega) \hat{f}(\omega)
    \end{aligned}
\end{equation}

\end{proof}

Figure~\ref{fig:atlas-box-convolution} interprets convolution as the length of overlap between two moving intervals.

\begin{figure}[!htbp]
\centering
\begin{tikzpicture}[>=Stealth]
\begin{axis}[width=5.8cm,height=4.8cm,axis lines=middle,ticklabel style={font=\small},label style={font=\small},legend style={font=\scriptsize,draw=black!20,fill=white},xlabel={$s$},xmin=-0.6,xmax=1.3,ymin=0,ymax=1.4,xtick={-0.25,0,0.75,1},xticklabels={$t-1$,$0$,$t$,$1$},ytick=\empty]
\addplot[figblue,very thick] coordinates {(0,1)(1,1)};
\addplot[figred,very thick] coordinates {(-0.25,0.65)(0.75,0.65)};
\addplot[figgreen,line width=4pt] coordinates {(0,0.3)(0.75,0.3)};
\draw[black!40,dashed] (axis cs:0,0) -- (axis cs:0,1.2);
\draw[black!40,dashed] (axis cs:0.75,0) -- (axis cs:0.75,1.2);
\end{axis}
\begin{axis}[width=5.8cm,height=4.8cm,axis lines=middle,ticklabel style={font=\small},label style={font=\small},legend style={font=\scriptsize,draw=black!20,fill=white},at={(6.2cm,0)},xlabel={$t$},ylabel={$(f*f)(t)$},xmin=-0.3,xmax=2.3,ymin=-0.1,ymax=1.4,xtick={0,1,2},ytick={1}]
\addplot[figblue,very thick] coordinates {(-0.3,0)(0,0)(1,1)(2,0)(2.3,0)};
\addplot[only marks,figgreen,mark=*,mark size=2.5pt] coordinates {(0.75,0.75)};
\draw[figgreen,dashed] (axis cs:0.75,0) -- (axis cs:0.75,0.75);
\end{axis}
\end{tikzpicture}
\caption{Left: at shift $t=3/4$, the intervals $[0,1]$ and $[t-1,t]$ overlap on a segment of length $3/4$. Right: as $t$ varies, that overlap length forms the triangular convolution of two unit-interval indicator functions. It is zero outside $[0,2]$ and reaches one at $t=1$.}
\label{fig:atlas-box-convolution}
\end{figure}

\subsection{The Heisenberg Uncertainty Principle}


\begin{theorem}[Heisenberg Uncertainty Principle]
For $f \in L^2(\R)$ with $\|f\|_{L^2}=1$, assuming $tf(t)$ and $\omega \hat{f}(\omega)$ are in $L^2(\R)$:
\begin{equation}
\left( \int_\R t^2 |f(t)|^2 \dd t \right) \left( \frac{1}{2\pi} \int_\R \omega^2 |\hat{f}(\omega)|^2 \dd \omega \right) \geq \frac{1}{4}
\end{equation}

\end{theorem}

\begin{proof}
Define the time and frequency second moments by
\[
\begin{gathered}
\Delta t^2 = \int_\R t^2 |f(t)|^2 \dd t \\
\Delta \omega^2 = \frac{1}{2\pi} \int_\R \omega^2 |\hat{f}(\omega)|^2 \dd \omega.
\end{gathered}
\]

For the derivation, we assume $f \in \mathcal{S}(\R)$ to ensure all boundary terms vanish and derivatives are well-behaved; the result extends to the specified domain in $L^2$ by a standard density argument.

\textit{Step 1: Relate $\Delta \omega^2$ to the time-domain derivative.}
First, we establish the derivative property of the Fourier transform. Using integration by parts on $\widehat{f'}(\omega)$:
\begin{equation}
\begin{aligned}
    \widehat{f'}(\omega)
    &= \int_\R f'(t) e^{-i\omega t} \dd t \\
    &= \left[ f(t) e^{-i\omega t} \right]_{-\infty}^\infty
    + \int_\R f(t) (i\omega) e^{-i\omega t} \dd t
    \end{aligned}
\end{equation}
Because $f \in \mathcal{S}(\R)$, $f(t) \to 0$ as $|t| \to \infty$, so the boundary term vanishes. Thus,
\[
\widehat{f'}(\omega) = i\omega \hat{f}(\omega).
\]
Applying Plancherel's theorem to the function $f'$:
\begin{equation}
\begin{aligned}
    \int_\R |f'(t)|^2 \dd t
    &= \frac{1}{2\pi} \int_\R |\widehat{f'}(\omega)|^2 \dd \omega \\
    &= \frac{1}{2\pi} \int_\R |i\omega \hat{f}(\omega)|^2 \dd \omega \\
    &= \frac{1}{2\pi} \int_\R \omega^2 |\hat{f}(\omega)|^2 \dd \omega
    = \Delta \omega^2
    \end{aligned}
\end{equation}

\textit{Step 2: Integration by parts on the time second moment.}
We evaluate the integral of $t (d/dt) |f(t)|^2$. Note that by the product rule:
\begin{equation}
\begin{aligned}
\frac{d}{dt} |f(t)|^2 &= \frac{d}{dt} (f(t) \overline{f(t)}) \\
&= f'(t) \overline{f(t)} + f(t) \overline{f'(t)} \\
&= 2 \Re\left( f'(t) \overline{f(t)} \right)
\end{aligned}
\end{equation}
Integrating by parts over $\R$:
\begin{equation}
\int_\R t \frac{d}{dt} |f(t)|^2 \dd t = \left[ t |f(t)|^2 \right]_{-\infty}^\infty - \int_\R |f(t)|^2 \dd t
\end{equation}

Since $f \in \mathcal{S}(\R)$, the boundary term vanishes. The remaining integral is $-\|f\|_{L^2}^2 = -1$. Equating the two expressions for the integral:
\begin{equation}
\begin{aligned}
-1 &= \int_\R t \left( f'(t) \overline{f(t)} + f(t) \overline{f'(t)} \right) \dd t \\
&= 2 \Re \int_\R t f'(t) \overline{f(t)} \dd t
\end{aligned}
\end{equation}

\textit{Step 3: Apply the Cauchy--Schwarz inequality.}
Taking the absolute value of both sides:
\begin{equation}
\begin{aligned}
1 &= \left| 2 \Re \int_\R t f'(t) \overline{f(t)} \dd t \right| \\
&\leq 2 \left| \int_\R t f'(t) \overline{f(t)} \dd t \right| \\
&\leq 2 \int_\R |t f(t)| |f'(t)| \dd t
\end{aligned}
\end{equation}

We apply the Cauchy--Schwarz inequality to the functions $A(t) = t f(t)$ and $B(t) = f'(t)$ in $L^2(\R)$:
\begin{equation}
\begin{aligned}
\int_\R |A(t) B(t)| \dd t &\leq \|A\|_{L^2} \|B\|_{L^2} \\
&= \left( \int_\R t^2 |f(t)|^2 \dd t \right)^{1/2} \left( \int_\R |f'(t)|^2 \dd t \right)^{1/2} \\
&= \sqrt{\Delta t^2 \Delta \omega^2}
\end{aligned}
\end{equation}
Substituting this bound back into the inequality:
\begin{equation}
1 \leq 2 \sqrt{\Delta t^2 \Delta \omega^2}
\end{equation}
Squaring both sides and dividing by 4 yields the uncertainty principle: $\Delta t^2 \Delta \omega^2 \geq 1/4$.

\textit{Step 4: Determine the equality condition.}
Equality in the Cauchy--Schwarz inequality holds if and only if $B(t) = c A(t)$ for some constant $c \in \C$, meaning $f'(t) = c t f(t)$.
Furthermore, equality in the bound $\Re(z) \leq |z|$ requires $z$ to be a nonnegative real number. Here,
\[
\begin{aligned}
z &= \int t f' \bar{f} \\
&= c \int t^2 |f|^2 \\
&= c \Delta t^2.
\end{aligned}
\]

Since $1 = -2 \Re(z)$, $z$ must be real and strictly negative, which implies $c$ is a negative real number. Let $c = -\alpha$ with $\alpha > 0$.
Solving the differential equation
\[
f'(t) = -\alpha t f(t):
\]

\begin{equation}
\frac{f'(t)}{f(t)} = -\alpha t \implies \ln f(t) = -\frac{\alpha}{2} t^2 + C \implies f(t) = A e^{-\frac{\alpha}{2} t^2}
\end{equation}
This explicitly proves that the lower bound is achieved if and only if $f$ is a Gaussian function.
\end{proof}

Figure~\ref{fig:ch14-uncertainty} illustrates the trade-off between localization in time and in frequency.

\begin{figure}[!htbp]
\centering
\begin{tikzpicture}
\begin{axis}[
  width=6.5cm,height=5.8cm, axis lines=middle,
  xlabel={$t$}, ylabel={$f(t)$},
  xmin=-3,xmax=3,ymin=0,ymax=1.3, xtick=\empty,ytick=\empty,
  ]
  \addplot[figblue,very thick,domain=-3:3,samples=150] {exp(-4*x^2)};
\end{axis}
\begin{axis}[
  at={(7.2cm,0)}, width=6.5cm,height=5.8cm, axis lines=middle,
  xlabel={$\xi$}, ylabel={$\hat f(\xi)$},
  xmin=-3,xmax=3,ymin=0,ymax=1.3, xtick=\empty,ytick=\empty,
  ]
  \addplot[figred,very thick,domain=-3:3,samples=150] {exp(-0.4*x^2)};
\end{axis}
\begin{axis}[
  at={(0,-5.2cm)}, width=6.5cm,height=5.8cm, axis lines=middle,
  xlabel={$t$}, ylabel={$f(t)$},
  xmin=-3,xmax=3,ymin=0,ymax=1.3, xtick=\empty,ytick=\empty,
  ]
  \addplot[figgreen,very thick,domain=-3:3,samples=150] {exp(-0.4*x^2)};
\end{axis}
\begin{axis}[
  at={(7.2cm,-5.2cm)}, width=6.5cm,height=5.8cm, axis lines=middle,
  xlabel={$\xi$}, ylabel={$\hat f(\xi)$},
  xmin=-3,xmax=3,ymin=0,ymax=1.3, xtick=\empty,ytick=\empty,
  ]
  \addplot[figblue,very thick,domain=-3:3,samples=150] {exp(-4*x^2)};
\end{axis}
\end{tikzpicture}
\caption[Uncertainty principle: time vs frequency]{Schematic time--frequency localization using Gaussian profiles. Left column: signals in time. Right column: frequency profiles. A narrow signal corresponds to a broad frequency profile (top row), while a broad signal corresponds to a narrow frequency profile (bottom row), illustrating the uncertainty bound discussed in the text.}
\label{fig:ch14-uncertainty}
\end{figure}

\section{Holomorphic Functions, Singularities, and Residues}
\label{sec:complex-analysis}

Complex differentiation and contour integration complement the preceding Fourier methods. They also connect the planar integral theorems to harmonic functions and provide tools for evaluating real integrals and Fourier transforms.

\subsection{Holomorphic Functions and the Cauchy--Riemann Equations}

\begin{definition}[Complex Differentiability and Holomorphy]
Let $U\subset\C$ be open. A function $f:U\to\C$ is \emph{complex differentiable} at $z_0\in U$ if the limit \[
f'(z_0):=\lim_{h\to0}\frac{f(z_0+h)-f(z_0)}{h}
\] exists independently of how the complex increment $h$ approaches zero. The function is \emph{holomorphic} on $U$ if it is complex differentiable at every point of $U$; it is holomorphic near $z_0$ if this holds in a neighborhood of $z_0$.
\end{definition}

This condition is stronger than real differentiability of a map $\R^2\to\R^2$. Writing
\[
\begin{gathered}
z=x+iy \\
f(z)=u(x,y)+iv(x,y).
\end{gathered}
\]
The difference quotient must have the same limit along the real and imaginary directions. This imposes relations between the partial derivatives of $u$ and $v$.

\begin{theorem}[Cauchy--Riemann Equations]
If $f=u+iv$ is complex differentiable at $z_0=x_0+iy_0$, then the partial derivatives of $u,v$ exist at $(x_0,y_0)$ and satisfy
\begin{equation}
\frac{\partial u}{\partial x}=\frac{\partial v}{\partial y},
\qquad
\frac{\partial u}{\partial y}=-\frac{\partial v}{\partial x}.
\end{equation}

\end{theorem}

\begin{proof}
\textit{Step 1: approach along the real direction.} Taking $h=\Delta x\in\R$ gives

\begin{align}
f'(z_0)
&=\lim_{\Delta x\to0}\left[
\frac{u(x_0+\Delta x,y_0)-u(x_0,y_0)}{\Delta x}\right.\notag\\
&\hspace{35mm}\left.+i\frac{v(x_0+\Delta x,y_0)-v(x_0,y_0)}{\Delta x}\right]\notag\\
&=u_x(x_0,y_0)+iv_x(x_0,y_0).
\end{align}

\textit{Step 2: approach along the imaginary direction.} Taking $h=i\Delta y$ with $\Delta y\in\R$ gives

\begin{align}
f'(z_0)
&=\lim_{\Delta y\to0}\left[
\frac{u(x_0,y_0+\Delta y)-u(x_0,y_0)}{i\Delta y}\right.\notag\\
&\hspace{35mm}\left.+\frac{v(x_0,y_0+\Delta y)-v(x_0,y_0)}{\Delta y}\right]\notag\\
&=-iu_y(x_0,y_0)+v_y(x_0,y_0),
\end{align}
because $1/i=-i$.

\textit{Step 3: compare the two limits.} Independence of the approach direction implies

\[
u_x+iv_x=v_y-iu_y.
\]
Equality of the real and imaginary parts gives $u_x=v_y$ and $v_x=-u_y$, as required.
\end{proof}

\begin{remark}[Sufficiency and Harmonic Functions]
If $u,v$ have continuous first partial derivatives on an open set and satisfy the Cauchy--Riemann equations there, then $f=u+iv$ is holomorphic. Moreover, the real and imaginary parts of a holomorphic function are harmonic. To see the differential mechanism, assuming $u,v\in C^2$, differentiate $u_x=v_y$ in $x$ and $u_y=-v_x$ in $y$:
\[
\begin{aligned}
\Delta u &= u_{xx}+u_{yy} \\
&= v_{yx}-v_{xy} \\
&= 0.
\end{aligned}
\]

Equality of mixed partial derivatives gives the cancellation; the argument for $\Delta v=0$ is analogous. Holomorphic functions in fact possess the required higher regularity, a result not proved here. This connects complex analysis to the Laplace equation and to the Dirichlet and Neumann problems discussed in Section~\ref{sec:pde-boundary-problems}.
\end{remark}

\subsection{Singularities, Residues, and the Residue Theorem}

\begin{definition}[Isolated Singularity and Residue]
An \emph{isolated singularity} of $f$ is a point $z_0$ around which $f$ is holomorphic on a punctured disk $0<|z-z_0|<r$, while its value at $z_0$ may be undefined or fail to give a holomorphic extension. On this punctured disk, $f$ has a \emph{Laurent expansion}
\begin{equation}
f(z)=\sum_{n=-\infty}^{\infty}c_n(z-z_0)^n.
\end{equation}
The coefficient of $(z-z_0)^{-1}$ is the \emph{residue}, denoted by
\[
\operatorname{Res}(f,z_0):=c_{-1}.
\]
The existence of the Laurent expansion is a standard result whose proof we omit.
\end{definition}

\begin{theorem}[Residue Theorem]
\label{thm:residue}
Let $\gamma$ be a positively oriented piecewise smooth simple closed curve. Suppose that $f$ is holomorphic on an open neighborhood of $\gamma$ and its interior, except for finitely many isolated singularities $z_1,\dots,z_k$ strictly inside $\gamma$. Then
\begin{equation}
\oint_\gamma f(z)\,\dd z
=2\pi i\sum_{j=1}^k\operatorname{Res}(f,z_j).
\end{equation}

\end{theorem}

\begin{proof}[Proof Sketch]
The argument rests on two facts. First, Cauchy's integral theorem states that the integral around a simple closed curve vanishes if $f$ is holomorphic on a neighborhood of the curve and its entire interior. Under $C^1$ regularity, Green's theorem (Theorem~\ref{thm:green-plane}) shows this directly: writing $f(z)\,\dd z=(u\,\dd x-v\,\dd y)+i(v\,\dd x+u\,\dd y)$, the two area integrands are $-v_x-u_y$ and $u_x-v_y$, both zero by the Cauchy--Riemann equations. Second, parametrizing a positively oriented circle by $z-z_0=\varepsilon e^{i\theta}$ gives
\begin{equation}
\oint_{|z-z_0|=\varepsilon}(z-z_0)^n\,\dd z
=\begin{cases}2\pi i,&n=-1,\\0,&n\in\Z,\ n\ne-1.\end{cases}
\end{equation}

Remove disjoint small disks around the singularities and apply Cauchy's theorem to the remaining region. The integral along $\gamma$ equals the sum of the counterclockwise integrals around the small circles. Integrating each Laurent expansion term by term selects only its coefficient $c_{-1}$, yielding the stated formula.
\end{proof}

\begin{example}[Evaluating a Real Integral by Residues]
Consider
\[
\int_{-\infty}^{\infty}\frac{\dd x}{1+x^2}.
\]

Its value is already accessible through the arctangent, but it illustrates a method that also applies when elementary antiderivatives are unavailable. The complex function $f(z)=1/(1+z^2)$ has simple poles at $z=\pm i$. For $R>1$, let $\gamma_R$ consist of the segment $[-R,R]$ followed by the upper semicircle
\[
C_R = \{Re^{i\theta}:0 \leq \theta \leq \pi\}.
\]
Only the pole at $i$ is enclosed, and
\begin{equation}
\operatorname{Res}(f,i)
=\lim_{z\to i}(z-i)f(z)
=\left.\frac{1}{z+i}\right|_{z=i}=\frac{1}{2i}.
\end{equation}
The residue theorem gives
\[
\oint_{\gamma_R}f(z)\,\dd z=\pi.
\]
On the semicircle, $|1+z^2|\geq R^2-1$, so
\begin{equation}
\left|\int_{C_R}f(z)\,\dd z\right|
\leq\frac{\pi R}{R^2-1}\longrightarrow0.
\end{equation}
Letting $R\to\infty$ therefore yields
\begin{equation}
\int_{-\infty}^{\infty}\frac{\dd x}{1+x^2}=\pi.
\end{equation}

Closing a real contour with an arc whose integral vanishes is also useful for explicit Fourier-transform calculations. For integrals containing $f(x)e^{i\omega x}$, the choice of half-plane must account for the exponential decay as well as the poles of $f$.
\end{example}

Figure~\ref{fig:atlas-residue-contour} identifies the enclosed pole and contour orientation in the residue calculation.

\begin{figure}[!htbp]
\centering
\begin{tikzpicture}[scale=1.15,>=Stealth]
\draw[->,black!40] (-3,0) -- (3.2,0) node[right] {$\mathrm{Re}$};
\draw[->,black!40] (0,-1.5) -- (0,3) node[above] {$\mathrm{Im}$};
\draw[figblue,very thick] (-2.5,0) -- (2.5,0) arc[start angle=0,end angle=180,radius=2.5];
\draw[figblue,very thick,->] (-1.8,0) -- (-0.8,0);
\draw[figblue,very thick,->] (35:2.5) arc[start angle=35,end angle=65,radius=2.5];
\node[below] at (-2.5,0) {$-R$};
\node[below] at (2.5,0) {$R$};
\node[above right,figblue] at (45:2.5) {$C_R$};
\node[figred] at (0,1) {$\times$};
\node[right,figred] at (0.1,1) {$i$};
\node[black!65] at (0,-1) {$\times$};
\node[right] at (0.1,-1) {$-i$};
\end{tikzpicture}
\caption{For $1/(1+z^2)$, the real segment from $-R$ to $R$ is closed by an upper semicircle traversed counterclockwise. Only the pole at $i$ lies inside; the pole at $-i$ lies outside. For $R>1$, the arc contribution tends to zero as $R$ grows, leaving the real-line integral.}
\label{fig:atlas-residue-contour}
\end{figure}